\chapter{现代人工智能方法提示}
\begin{enumerate}[leftmargin=*]
  \item shape 是语义轴，dtype/device 决定数值算子，梯度状态决定是否进入反向图。
  \item 比较局部感受野、递归状态和内容寻址；所有路径都要服从因果可见集。
  \item GNN 增加关系图；DRL 增加状态、动作、转移和长期奖励。
  \item 依次写仿射、激活、仿射和链式法则；残差提供恒等梯度路径。
  \item 输出是相邻差；检查填充是否让 $t$ 读取 $t+1$。
  \item 同步置换 $Q,K,V$ 只会重排输出；手算 $\operatorname{softmax}(0,1)$。
  \item 图边要有生效时间；RL 转移模型必须响应动作。
  \item 均值应不变，顺序模型应变化；两类 mask 分别处理未来和不存在位置。
  \item checkpoint 同时绑定配置、参数和输入 schema。
  \item 为每个决策日截断边表，而不是今天的全图。
  \item 查是否有可训练序列层、位置结构和置换负例。
  \item 比较可训练参数和基础表示是否更新；记录权重训练截止与语料版本。
  \item 先用两步、有限状态的动态规划给出解析策略。
  \item 项目只选一条主线，但必须保留树或线性基线并连接净收益。
\end{enumerate}
